\chapter{Gaiotto--Rapcak $Y$-algebras and the triality $S_3$}
\label{ch:y-algebras-triality}

Gaiotto--Rapcak's three-corner vertex algebras $Y_{L,M,N}^\Psi$ arise as
boundary algebras of intersecting M5-brane corners in $\mathbb{C}^3$ with
$\Omega$-background \cite{GaiottoRapcak2017}. Their defining feature is
a \emph{triality} groupoid with symmetric group $S_3$ acting on the
generic parameter space, corresponding to the Weyl symmetry of the three
$\Omega$-background parameters $(q_1, q_2, q_3)$ subject to
$q_1 q_2 q_3 = 1$. We construct the induced $S_3$ action on each slot of
the ordered-bar fingerprint $\varphi'(A)$
(Chapter~\ref{chap:infinite-fingerprint}), prove equivariance of the
chiral quantum group datum $\sQ_g^{k,f,\mu}$ from
Theorem~\ref{thm:unified-chiral-QG}, and identify the induced finite
symmetry of the Koszulness moduli scheme
$M_{\mathrm{Kosz}}$ (Chapter~\ref{chap:koszulness-moduli}).

Throughout we work on the generic locus of the universal
$Y_{L,M,N}^{\Psi}$ family on which the Gaiotto--Rapcak triality
isomorphisms are defined. Diagonal stabilisers and special
$\Psi$-values are treated at the end.

\section{Setup: $Y_{L,M,N}^\Psi$ from three-corner M5 branes}
\label{sec:y-algebra-setup}

\begin{definition}[Gaiotto--Rapcak $Y$-algebra; {\cite[Def.~2.1]{GaiottoRapcak2017}}]
\label{def:y-algebra}
Fix non-negative integers $(L,M,N) \in \bZ_{\ge 0}^3$ and a
\emph{level parameter} $\Psi \in \bC \setminus \{0\}$. The
Gaiotto--Rapcak $Y$-algebra $Y_{L,M,N}^\Psi$ is the vertex algebra
obtained as the BRST reduction of three corners of
$\mathcal{N}=4$ super-Yang--Mills with gauge groups
$(\mathrm{U}(L), \mathrm{U}(M), \mathrm{U}(N))$ meeting at a single
point in $\bR \times \bC$, under the holomorphic-topological twist of
\cite{CostelloGaiotto2018}. Equivalently, $Y_{L,M,N}^\Psi$ is a
quantum Drinfeld--Sokolov reduction of
$\widehat{\mathfrak{gl}}_{L+M+N}$ at level
$k_\Psi(L,M,N)$ determined by $\Psi$ and the three-box Young diagram
shape $(L,M,N)$.
\end{definition}

The $\Omega$-background parameters $(\epsilon_1, \epsilon_2, \epsilon_3)$
with $\epsilon_1 + \epsilon_2 + \epsilon_3 = 0$ multiplicatively
re-encode as $(q_1, q_2, q_3) = (e^{\epsilon_1}, e^{\epsilon_2}, e^{\epsilon_3})$
with $q_1 q_2 q_3 = 1$. The label correspondence is $(L,M,N) \leftrightarrow
(\epsilon_1, \epsilon_2, \epsilon_3)$ up to a rescaling depending on
$\Psi$:
\begin{equation}
\label{eq:y-algebra-level}
\epsilon_i = \frac{\sigma_i(L,M,N; \Psi)}{\Psi - 1}, \qquad i = 1,2,3,
\end{equation}
where $\sigma_i$ are the three corner-length functions in the
Gaiotto--Rapcak parametrisation (see
\cite[Eq.~3.8]{GaiottoRapcak2017}).

\begin{definition}[Triality action]
\label{def:triality-S3}
The \emph{triality} action is the Gaiotto--Rapcak--Creutzig--Linshaw
$S_3$-action on the generic parameter space:
\[
  \sigma\cdot(L,M,N,\Psi)=(L_\sigma,M_\sigma,N_\sigma,\Psi_\sigma),
  \qquad \sigma\in S_3 .
\]
It permutes the three corners and applies the corresponding birational
change of level parameter, equivalently the Weyl permutation
$(q_1,q_2,q_3)\mapsto\sigma(q_1,q_2,q_3)$ of the $\Omega$-background
torus. We write
\[
  Y_{\sigma\cdot(L,M,N,\Psi)}
  :=
  Y_{L_\sigma,M_\sigma,N_\sigma}^{\Psi_\sigma}.
\]
\end{definition}

\begin{remark}[$S_3$ is algebra-specific, not operadic]
\label{rem:triality-scope}
The triality $S_3$ is \emph{algebra-specific}: it arises from the
three-corner symmetry of the specific construction
$Y_{L,M,N}^\Psi$ via its M5-brane origin. It is \emph{not} an
operadic $S_3$ on $E_3$; the counterexample $k[x]/(x^2)$ is
$E_3$ as a dg algebra but has no triality. The Miki $S_3$ on the
quantum toroidal algebra $U_{q,t}(\widehat{\widehat{\mathfrak{gl}}}_1)$
is the image of our $S_3$ under the Miura transform. We never derive
triality from $E_3$ data alone; we construct it from the specific
$(L,M,N)$ parametrisation.

$Y_{L,M,N}^\Psi$ is \emph{natively} $E_1$-chiral; the $E_2$
and $E_3$ structures live on its derived centre. All triality claims
below are stated on the native object and transported to its derived
centre through the constructions of
Chapter~\ref{ch:dnp-identification-master} and
Chapter~\ref{ch:chiral-higher-deligne}.

Every appearance of ``triality preserves $X$'' below is backed by an
explicit arrow $\sigma: X(Y_{L,M,N}^\Psi) \xrightarrow{\sim}
X(Y_{\sigma\cdot(L,M,N,\Psi)})$; the arrow is the content of the
preservation claim.
\end{remark}

\section{Explicit $S_3$ on Creutzig--Linshaw strong generators}
\label{sec:triality-strong-generators}

The first step is to exhibit triality as a permutation of
\emph{strong generators} of $Y_{L,M,N}^\Psi$, compatibly with the PBW
filtration. We follow the Creutzig--Linshaw presentation
\cite{CreutzigLinshaw2019}.

\begin{proposition}[Creutzig--Linshaw strong generators under $S_3$; {\cite{CreutzigLinshaw2019}}]
\label{prop:strong-generators-s3}
There is a family of strong generators $W^{(n)}_{\sigma}$ for
$n \ge 2$ and $\sigma \in S_3$ indexing the three corners of
$Y_{L,M,N}^\Psi$, such that:
\begin{enumerate}
\item[(i)] $W^{(n)}_{\sigma}$ is a Virasoro primary of conformal
 weight $n$ on corner $\sigma$.
\item[(ii)] The $S_3$ action on labels induces
 $\sigma_* W^{(n)}_{\tau} = W^{(n)}_{\sigma\tau}$, permuting generators
 across corners without changing conformal weight or PBW degree.
\item[(iii)] The PBW filtration $F^\bullet Y_{L,M,N}^\Psi$ is
 $S_3$-equivariant: $\sigma_* F^p = F^p$.
\end{enumerate}
\end{proposition}

\begin{proof}
Part (i) is \cite[Thm.~3.4]{CreutzigLinshaw2019} for the
$W$-algebra corner structure; part (ii) follows from the manifest
three-corner symmetry of the Coulomb-branch realisation
\cite[\S4]{ProchazkRapcak2018} where the $S_3$ acts by permuting
the three generating screening charges; part (iii) is then
automatic because the PBW filtration is determined by the
conformal-weight filtration
$F^p Y = \bigoplus_{n \le p} Y[n]$ (equivalently, strong-generator
Li filtration), and $S_3$ preserves conformal weight by (i).
\end{proof}

\begin{remark}[explicit $S_3$ on $W_\infty$ generators]
\label{rem:wn-explicit-triality}
For the pure-$\mathcal{W}_\infty$ corner (sending $L$ or $M$ or $N$
to zero), the $S_3$ action restricts to the known triality of
$\mathcal{W}_\infty[\lambda]$ at the three values
$\lambda \in \{L, M, N\}$ (Gaberdiel--Gopakumar triality
\cite{GaberdielGopakumar2012}), permuting the three dual free-field
realisations. This is the classical ``Gaberdiel--Gopakumar $S_3$'';
the Gaiotto--Rapcak $S_3$ on $Y_{L,M,N}^\Psi$ reduces to this at the
boundary.
\end{remark}

\section{Triality preserves the fingerprint}
\label{sec:triality-fingerprint}

Chapter~\ref{chap:infinite-fingerprint} defines the canonical
ordered-bar fingerprint
\[
\varphi'(A) =
(p_{\max}, r_{\max}, \chi_{\mathrm{VOA}},
n_{\mathrm{strong}}, \mathrm{coset}, \kappaChHodge)
\]
for a chirally Koszul algebra $A$. The main theorem of this chapter is:

\begin{theorem}[Triality preserves the fingerprint]
\label{thm:triality-preserves-fingerprint}
For all $\sigma \in S_3$ and all generic $(L,M,N,\Psi)$,
\[
\varphi'\bigl(Y_{\sigma\cdot(L,M,N,\Psi)}\bigr)
\;=\;
\varphi'\bigl(Y_{L,M,N}^\Psi\bigr).
\]
In particular, triality preserves every coordinate of the
fingerprint:
(a) pole order $p_{\max}$; (b) shadow depth $r_{\max}$ (hence
shadow class $\in \{G, L, C, M, FF\}$); (c) VOA character
$\chi_{\mathrm{VOA}}$; (d) minimal strong-generator count
$n_{\mathrm{strong}}$; (e) commutant coset class;
(f) $\kappaChHodge$.
\end{theorem}

\begin{proof}
We prove each slot separately. Throughout, let
$A = Y_{L,M,N}^\Psi$ and
$A^\sigma = Y_{\sigma\cdot(L,M,N,\Psi)}$;
by Proposition~\ref{prop:strong-generators-s3} there is an
$S_3$-equivariant isomorphism of the underlying graded vector spaces
$\sigma_*: A \to A^\sigma$ preserving conformal weight and PBW
degree.

\smallskip
\emph{(a) Pole order $p_{\max}$.} Pole order is the maximum
$p$ such that the OPE $W^{(m)}_\rho(z) W^{(n)}_\tau(w) \sim
C_{\rho\tau}^{(p)} (z-w)^{-p} + \dotsb$ has a nonzero coefficient, as
$(m, n, \rho, \tau)$ range over strong generators. Because
$\sigma_*$ is a vertex algebra isomorphism (Proposition
\ref{prop:strong-generators-s3}\,(i)-(iii) plus OPE-preservation, which
follows from the three-corner symmetry of the Coulomb-branch BRST
reduction \cite[\S5]{ProchazkRapcak2018}), OPE coefficients are
permuted bijectively. Hence the set of pole orders is
$S_3$-invariant, so is its maximum: $p_{\max}(A^\sigma) = p_{\max}(A)$.

\smallskip
\emph{(b) Shadow depth $r_{\max}$.} The shadow tower
$\{S_r\}_{r \ge 2}$ is extracted from the bar complex
$B^{\mathrm{ord}}(A) = T^c(s^{-1}\bar A)$ with deconcatenation
coproduct, via the Volume~I recursion
$S_r = \operatorname{Res}^{\mathrm{coll}}_{0,r}(\Theta_A)$
(Vol.~I, Theorem~\ref*{V1-thm:mc2-bar-intrinsic}). Since the bar
complex depends \emph{only} on the augmented vertex algebra
structure (OPE data + vacuum + translation) and $\sigma_*$ is an
OPE-preserving vertex algebra isomorphism by (a), we have
$B^{\mathrm{ord}}(A^\sigma) \cong \sigma_* B^{\mathrm{ord}}(A)$
as dg coalgebras. The shadow tower is therefore mapped isomorphically
by $\sigma_*$, yielding $S_r(A^\sigma) = \sigma_* S_r(A)$. In
particular nonvanishing is preserved, so
$r_{\max}(A^\sigma) = r_{\max}(A)$, and the shadow class
(G / L / C / M / FF) is preserved.

\smallskip
\emph{(c) VOA character.} The graded character
$\chi_{\mathrm{VOA}}(q) = \tr_A q^{L_0 - c/24}$
is determined by the conformal vector and the $L_0$-grading, which are
preserved by the triality isomorphism
(Proposition~\ref{prop:strong-generators-s3}\,(i) and
\cite{CreutzigLinshaw2019}). Hence the central charge is invariant
and $\chi_{\mathrm{VOA}}(q)$ is preserved.

\smallskip
\emph{(d) Minimal strong-generator count $n_{\mathrm{strong}}$.}
Proposition~\ref{prop:strong-generators-s3}\,(ii) shows that
$\sigma_*$ bijects the strong-generator set; minimality is a
cardinality condition, so is preserved.

\smallskip
\emph{(e) Coset class.} The commutant pair
$(\operatorname{Com}(\mathcal{H}_\kappa, A), A/\operatorname{Com})$
depends only on the bar-intrinsic $\kappaChHodge(A)$ and the
Heisenberg subalgebra structure, both $S_3$-invariant by (a) and (c).
Hence $\sigma$ induces an isomorphism on coset pairs, preserving
the coset class in
$\pi_0(\mathbf{Coset})$.

\smallskip
\emph{(f) Chiral curvature $\kappaChHodge$.}
On the affine branch $\{L=0\}$ and its triality images,
$\kappaChHodge(Y)=c(Y)/2$ by the Vol.~I landscape formula for
affine chiral algebras. On the interior, the Drinfeld--Sokolov
construction transports this scalar through the same conformal vector
and BRST complex. Since triality is an isomorphism of conformal vertex
algebras, the central charge and the transported DS scalar agree on
$A$ and $A^\sigma$.

\smallskip
Combining (a)-(f) with
the product decomposition
\[
\varphi'(A) = (p_{\max}, r_{\max}, \chi_{\mathrm{VOA}},
n_{\mathrm{strong}}, \mathrm{coset}, \kappaChHodge)
\]
gives the theorem.
\end{proof}

\begin{corollary}[Triality preservation of shadow class]
\label{cor:triality-shadow-class}
The Gaiotto--Rapcak triality preserves shadow class for all
generic $(L,M,N,\Psi)$. This is exactly
Theorem~\ref{thm:triality-preserves-fingerprint}\,(b).
\end{corollary}

\section{Triality equivariance of $\kappaChHodge$}
\label{sec:triality-kappa}

\begin{theorem}[Triality preserves $\kappaChHodge$]
\label{thm:triality-preserves-kappa}
For all $\sigma \in S_3$ and all generic $(L,M,N,\Psi)$,
\[
\kappaChHodge\bigl(Y_{\sigma\cdot(L,M,N,\Psi)}\bigr)
\;=\;
\kappaChHodge\bigl(Y_{L,M,N}^\Psi\bigr).
\]
\end{theorem}

\begin{proof}
This is the sixth coordinate of
Theorem~\ref{thm:triality-preserves-fingerprint}.
\end{proof}

\begin{remark}[Subscript convention on $\kappa$]
\label{rem:kappa-subscript-y-algebra}
Every $\kappa$ in this chapter carries an explicit subscript; we
write $\kappaChHodge$ for the chiral modular characteristic.
The categorical $\kappa_{\mathrm{cat}}$, BKM $\kappa_{\mathrm{BKM}}$,
and fiber $\kappa_{\mathrm{fiber}}$ are not applicable to the ambient
$E_1$-chiral $Y_{L,M,N}^\Psi$; they arise only after CY or BKM
enhancements (Volume~III Chapter~\texttt{V3-ch:cy-to-chiral-functor}).
\end{remark}

\section{Triality lifts to the chiral QG datum}
\label{sec:triality-chiral-qg}

The unified chiral quantum group theorem
(Theorem~\ref{thm:unified-chiral-QG}) constructs for each chiral
Koszul algebra $A$ in the standard landscape a datum
\(
\sQ(A) \;=\; \bigl(A,\, \Delta_z,\, R(z),\, \epsilon,\, S\bigr).
\)
For $Y$-algebras the datum is
$\sQ(Y_{L,M,N}^\Psi)$. Triality lifts:

\begin{theorem}[$\sQ$ is $S_3$-equivariant]
\label{thm:triality-lifts-to-chiral-qg}
The assignment $(L,M,N,\Psi) \mapsto \sQ(Y_{L,M,N}^\Psi)$ is
$S_3$-equivariant: for every $\sigma \in S_3$, the isomorphism
$\sigma_*: Y_{L,M,N}^\Psi \xrightarrow{\sim}
Y_{\sigma\cdot(L,M,N,\Psi)}$
from Proposition~\ref{prop:strong-generators-s3} extends uniquely
to an isomorphism of chiral QG data,
\[
\sigma_*:\; \sQ\bigl(Y_{L,M,N}^\Psi\bigr)
\;\xrightarrow{\sim}\;
\sQ\bigl(Y_{\sigma\cdot(L,M,N,\Psi)}\bigr),
\]
preserving spectral coproduct $\Delta_z$, spectral $R$-matrix
$R(z)$, counit $\epsilon$, and antipode $S$.
\end{theorem}

\begin{proof}
The four pieces of the datum admit explicit formulas in terms of the
vertex algebra $Y_{L,M,N}^\Psi$ and its Koszul dual bar complex;
each formula is manifestly $S_3$-equivariant. We argue the four
pieces in turn.

\smallskip
\emph{$\Delta_z$.} The Drinfeld spectral coproduct is built from
the twisting morphism
$\pi: B^{\mathrm{ord}}(A) \to A$ and the Koszul-dual bar
co-deconcatenation; both are functorial in VA isomorphisms and
commute with $\sigma_*$ by Theorem~\ref{thm:triality-preserves-fingerprint}\,(a,b,d).

\smallskip
\emph{$R(z)$.} The spectral $R$-matrix is
$R(z) = \exp(r(z)) \bmod \hbar^2$-corrections, where
$r(z) = \operatorname{Res}^{\mathrm{coll}}_{0,2}(\Theta_A)$ is the
binary collision residue (Chapter~\ref{ch:dnp-identification-master}).
Theorem~\ref{thm:triality-preserves-fingerprint}\,(b) gives
$\sigma_* r(z) = r(z)$ for the $S_3$-image, so $R(z)$ is preserved.

\smallskip
\emph{$\epsilon$.} The counit is the vacuum projector
$\epsilon: A \twoheadrightarrow \bC$, which is $S_3$-canonical
(the vacuum $\ket{0}$ is $S_3$-invariant as it lies in the
weight-zero subspace $A[0] = \bC \ket{0}$).

\smallskip
\emph{$S$.} The antipode is constructed from the Shapovalov form
$\langle -, - \rangle$ via $S = -\text{flip}$; Shapovalov is built
from the OPE and the $L_0$-grading, both $S_3$-equivariant, hence
$S$ is $S_3$-equivariant.

The four equivariances combine to an isomorphism of quintuples,
proving the theorem.
\end{proof}

\begin{corollary}[triality in the unified $\sQ_g^{k,f,\mu}$]
\label{cor:triality-unified-Q}
Under the unified theorem identifying
$\sQ(Y_{L,M,N}^\Psi) = \sQ_{\mathfrak{gl}_{L+M+N}}^{k_\Psi, f_{L,M,N}, 0}$
with principal three-block nilpotent $f_{L,M,N}$, the $S_3$-action
on $(L,M,N)$ induces the expected permutation action on block
shapes, matching the outer automorphism of the Weyl group of
$\mathfrak{gl}_{L+M+N}$ preserving the three-block partition.
\end{corollary}

\section{Descent to Koszulness moduli and GRT$_1$}
\label{sec:triality-koszulness-moduli}

Chapter~\ref{chap:koszulness-moduli} constructs a
GRT$_1$-equivariant moduli scheme
$M_{\mathrm{Kosz}}(A)$ of Koszulness characterisations. Triality
acts on this scheme by functorial transport of Koszul-duality data.

\begin{corollary}[triality normalises the Koszulness moduli action]
\label{cor:triality-normalizes-mkosz}
For every generic $Y_{L,M,N}^{\Psi}$ and every $\sigma\in S_3$, the
triality isomorphism
\[
\sigma_*:Y_{L,M,N}^{\Psi}\xrightarrow{\sim}
Y_{\sigma\cdot(L,M,N,\Psi)}
\]
induces an isomorphism of GRT$_1$-torsors
\[
\sigma_*:
M_{\mathrm{Kosz}}\!\bigl(Y_{L,M,N}^{\Psi}\bigr)
\xrightarrow{\sim}
M_{\mathrm{Kosz}}\!\bigl(Y_{\sigma\cdot(L,M,N,\Psi)}\bigr).
\]
On a stabiliser stratum, this gives an invariant-subscheme action of the
finite stabiliser $S_3(L,M,N,\Psi)\subset S_3$ on
$M_{\mathrm{Kosz}}(Y_{L,M,N}^{\Psi})$. This finite symmetry is external
to $\mathrm{GRT}_1$: since $\mathrm{GRT}_1$ is prounipotent over
$\mathbb Q$, every finite subgroup of $\mathrm{GRT}_1(\mathbb Q)$ is
trivial.
\end{corollary}

\begin{proof}
A point of $M_{\mathrm{Kosz}}(A)$ is a Koszul-duality datum
$(A,A^!,\Phi)$ with an associator coordinate. Applying $\sigma_*$ to
the algebra and the bar complex gives
\[
(A,A^!,\Phi)\longmapsto
(A^\sigma,(A^\sigma)^!,\Phi).
\]
The associator coordinate is unchanged, so this transport commutes with
the GRT$_1$ action on associators. If $\sigma$ fixes the parameter, the
same construction is an automorphism of the fibre, and the stabiliser
acts on the corresponding fixed subscheme. Finally, a finite subgroup
of a prounipotent group over a characteristic-zero field is trivial:
after embedding into a unipotent matrix group, every finite-order
unipotent element has all eigenvalues equal to $1$ and logarithm zero,
hence is the identity.
\end{proof}

\section{Caveat: degenerate loci}
\label{sec:triality-caveat}

\begin{theorem}[accidental enhancements at degenerate loci]
\label{thm:triality-obstruction-non-generic}
Let $\Delta \subset \bZ_{\ge 0}^3 \times \bC$ be the union of the
three diagonals $\{L = M\}, \{M = N\}, \{L = N\}$ and the poles or
cusps of the level parameter. Then:
\begin{enumerate}
\item[(i)] For $(L,M,N,\Psi) \notin \Delta$ (generic locus), the
 $S_3$-action on $\sQ(Y_{L,M,N}^\Psi)$ is \emph{free}: no
 nontrivial $\sigma$ fixes $(L,M,N,\Psi)$, and the orbit size is $6$.
\item[(ii)] On the diagonal $\{L=M\}$ and its $S_3$-images, the
 stabiliser is the involution $(L\,M)$, and the action has orbit
 size $3$. The induced action on $\sQ$ is by an
 involutive vertex-algebra automorphism with nontrivial fixed
 subalgebra; the quotient is a $\bZ/2$-orbifold of
 $Y_{L,L,N}^\Psi$.
\item[(iii)] At $L = M = N$, the stabiliser is the full $S_3$,
 orbit size $1$. The induced action on $\sQ$ is through an
 outer $S_3$-automorphism; the invariant subalgebra
 $\bigl(Y_{L,L,L}^\Psi\bigr)^{S_3}$ is a diagonal coset and is a
 proper subalgebra.
\item[(iv)] At a pole or cusp of $\Psi$ the algebra degenerates; the
 fingerprint still coincides across $S_3$-orbits (by
 Theorem~\ref{thm:triality-preserves-fingerprint}, which holds at
 every parameter where $Y_{L,M,N}^\Psi$ is defined), and the
 chiral QG datum $\sQ$ acquires the extra automorphisms of the
 limiting free-field or infinite-level model.
\end{enumerate}
\end{theorem}

\begin{proof}
(i) is free-ness of the permutation action on pairwise-distinct
triples. (ii) follows from Remark~\ref{rem:wn-explicit-triality}
and the explicit fixed-point computation on
strong generators: $(L\,M)$ fixes exactly the diagonal generators
$W^{(n)}_{\mathrm{diag}} = (W^{(n)}_L + W^{(n)}_M)/\sqrt{2}$, giving the
$\bZ/2$-orbifold. (iii) follows by iteration: the three
involutions generate $S_3$, and the joint fixed subalgebra is the
invariant-polynomial subalgebra under the corner-permutation action
on three equal-size Young tableaux; the quotient is the corresponding
diagonal coset (a known construction
\cite{ProchazkRapcak2018}).
(iv) follows from the standard free-field and infinite-level
descriptions of the corner construction: in either limit the screening
complex drops rank, so the automorphism group contains the corner
permutation group and the automorphisms of the limiting free fields.
\end{proof}

\section{Consequences}
\label{sec:triality-summary}

Triality preserves the ordered-bar fingerprint on the generic locus,
and the stabiliser stratification is governed by
Theorem~\ref{thm:triality-obstruction-non-generic}. The same arrows
give the following uses.
\begin{itemize}
\item Chapter~\ref{ch:dnp-identification-master} uses
 Theorem~\ref{thm:triality-preserves-fingerprint}\,(b) to
 $S_3$-orbit the seven-face star.
\item Chapter~\ref{ch:chiral-higher-deligne} uses
 Theorem~\ref{thm:triality-lifts-to-chiral-qg} to lift triality
 to $Z^{\mathrm{der}}_{\mathrm{ch}}(Y_{L,M,N}^\Psi)$ as an
 $S_3$-action on the $E_3$-topological bulk.
\item Chapter~\ref{chap:koszulness-moduli}'s
 $\mathrm{GRT}_1$-torsor is normalised by the external triality
 symmetry in Corollary~\ref{cor:triality-normalizes-mkosz}.
\end{itemize}
